% =============================================================================
% Template dissertation using the qdissertation class
% =============================================================================
% This file is the main entry point. Build: make (uses latexmk; see Makefile).
% Default engine: lualatex. To use another: make ENGINE=pdf or make ENGINE=xe.
%
% SPDX-FileCopyrightText: 2026 Frank Langbein <frank@langbein.org>
% SPDX-License-Identifier: LPPL-1.3c
% SPDX-License-Identifier: CC-BY-SA-4.0
%
% This template and other .tex/.bib files in this directory are part of the
% same work and may be distributed and/or modified under the LaTeX Project
% Public License (LPPL), version 1.3c. See LICENSE file.
%
% -----------------------------------------------------------------------------
% Overview (qdissertation class)
% -----------------------------------------------------------------------------
% The class is a wrapper around the standard book class. Layout: A4, 12pt,
% 3cm margins; single line spacing; ragged-right; main text 12pt, footnotes and
% captions at least 11pt; preliminary pages in roman numerals, then Arabic
% for main matter. Chapter style: top bar in header, ``Chapter N'' or
% ``Appendix A'' in header; number/letter and title on first line of chapter.
%
% -----------------------------------------------------------------------------
% Report structure
% -----------------------------------------------------------------------------
% Main body: Introduction, Background, Methodology (or ``Specification, Design
% and Implementation'' / ``Approach''), Results and Evaluation, Conclusions and
% Future Work, Reflection on Learning. Chapter~3 can be adapted or split; not all
% sections apply to every project (system design, research, AI/ML). See each
% chapter file for details and alternatives.
%
% Suggested sections at a glance (omit/rename as needed):
%   1 Introduction:    Motivation and context; Aims and objectives; Contributions;
%                      Dissertation structure. Optional: Scope, Terminology.
%   2 Background:      Context; Related work (or Literature review); Problem statement.
%                      Optional: Theory and concepts; Stakeholders and constraints.
%   3 Methodology:     Depends on project type. System design: Specification, System
%                      design, Implementation. Research: Research design, Literature
%                      and evidence, Approach and solution, Implementation. AI/ML:
%                      Specification, Datasets, AI/ML models, System design, Implementation.
%                      Optional: Ethical considerations; Summary.
%   4 Results:         Setup (or Experimental setup); Results; Discussion; Limitations.
%                      Optional: Ablation studies; Comparison with baselines.
%   5 Conclusions:     Conclusions; Future work.
%   6 Reflection:      Critical narrative; Analysis and application; Implications
%                      for future practice. (Omit chapter if not required.)
%
% -----------------------------------------------------------------------------
% Variations (degree, structure)
% -----------------------------------------------------------------------------
% Select your degree via \degree below; commented alternatives shown for
% BSc, MSc.
%
% -----------------------------------------------------------------------------
% Class options
% -----------------------------------------------------------------------------
% Font: mixed (default) = sans text + serif math; sans = all sans-serif;
%   serif = serif text and math; hacker = typewriter/monospace open-source font throughout.
% Layout: oneside (default) or twoside.
% Citation: numeric (default) = plainnat, [1], [1,2]; harvard = agsm, (Author, 2000).
% Alignment: raggedright (default) = ragged-right; justified = fully justified text.
% Line spacing: single (default); onehalfspacing; or linespread115/linespread12
%   for institutions allowing 1.0--1.3.
% Paragraph style: parskip (default) = space between paragraphs, no first-line indent;
%   parindent = traditional first-line indent, minimal parskip.
% Optional typography:
%   fullwidth  = enable \begin{fullwidth}...\end{fullwidth} and fullwidthfigure/
%                fullwidthtable environments that extend into the margin (requires
%                changepage or chngpage).
%   italic = chapter, section, subsection and captions in italic (default: bold/roman).
% Engine: pdfTeX vs XeLaTeX/LuaLaTeX is detected for font setup.
%
% Defaults (no option): link colour DarkBlue, cite/URL colour DarkGreen; tight list
% margins; quote, quotation, and verse in smaller type with indents;
% \epigraph{quote}{source} for chapter/section-opening quotes.
% See Section~\ref{sec:ch1-reference} in Chapter~1 for examples of \newthought,
% fullwidth, fullwidthfigure, quote, and epigraph.
%
% -----------------------------------------------------------------------------
% Metadata and title page
% -----------------------------------------------------------------------------
% \title, \author, \degree, \university, \date drive the generated title page.
%
% -----------------------------------------------------------------------------
% License and copyright (SPDX)
% -----------------------------------------------------------------------------
% \copyrightholder defaults to \author; \licenseyear defaults to current year.
% When \license is set: (1) a copyright page is inserted after the abstract;
% (2) a License chapter is inserted just before the bibliography.
% If qdissertation-licenses.tex is present, the class loads it for full
% license text; otherwise name + URL only.
%
% -----------------------------------------------------------------------------
% Float placement (figures, tables, algorithms)
% -----------------------------------------------------------------------------
% Placement letters: h = here, t = top, b = bottom, p = float page. E.g.\@
% [htbp], [t], [!t]. Use [H] (float pkg) sparingly. The class relaxes float
% parameters so more floats fit per page.
%
% -----------------------------------------------------------------------------
% Headings: use title case for chapter titles (capitalise major words; keep
% articles, conjunctions, and short prepositions lowercase unless first/last
% word). Use sentence case for sections, subsections, and below (only the
% first letter and proper nouns capitalised). Or adjust as preferred, but keep
% it consistent.
%
% -----------------------------------------------------------------------------
% LaTeX formatting: \citep/\citet (not \cite); ~ before \citep/\citet after a word;
% \@ after abbreviation; ~ for non-breaking (Fig.~1); -- en-dash, --- em-dash.
%
% -----------------------------------------------------------------------------
% Reviewer comments: %ID (e.g.\@ %F) at line start, space after ID. Include
% comments directly in the text and discussions or questions can be included
% as comments afterwards. Generally, these comments should be inserted closely
% after the sentence or element it refers to. Delete the commetns when
% addressed; do not keep ``for checking''. Edits can be easier found via git
% diffs. Find comments with: grep -r '^%[A-Za-z] ' *.tex */*.tex
%
% -----------------------------------------------------------------------------
% Footnotes: prefer putting short asides in parentheses in the main text. Use
% footnotes only sparingly when they add value (e.g.\@ translations, optional
% clarifications that would interrupt the flow).
%
% -----------------------------------------------------------------------------
% Word count: run ``make wordcount'' (texcount -inc; per-file and total with
% caption words shown separately so you can apply your institution's rules).
% Word limits typically apply to main body only; exclude front matter, captions,
% appendices. Confirm your institution's word limits and structure before submission.
%
% -----------------------------------------------------------------------------
% Submission: ensure PDF has embedded fonts (pdflatex does this by default),
% and is compatible with your institution's requirements (e.g. no encryption, be
% careful with obfuscation).
% =============================================================================
\documentclass[mixed,oneside,numeric,raggedright,fullwidth]{qdissertation}

% -----------------------------------------------------------------------------
% Optional: use a different acronyms file
% -----------------------------------------------------------------------------
%\renewcommand{\qdissertationacronymfile}{my-acronyms.tex}

% -----------------------------------------------------------------------------
% Useful packages for dissertation writing
% -----------------------------------------------------------------------------

% Tables
\usepackage{tabularray}      % Flexible column types, rules, longtblr for multi-page tables
\usepackage{booktabs}        % Alternative (traditional): \toprule, \midrule, \bottomrule
\usepackage{longtable}       % Tables that span multiple pages (or use tabularray longtblr)
%\usepackage{multirow}        % Span multiple rows in tables
%\usepackage{multicol}        % Multiple columns in text/tables
%\usepackage{colortbl}        % Coloured rows/cells in tables

% Math, symbols, related
\usepackage{mathtools}       % Extensions to amsmath (matrix envs, \DeclarePairedDelimiter)
\usepackage{pifont}          % Dingbat symbols, e.g. \ding{51}, \ding{43}
\usepackage{siunitx}         % Units and numbers: \SI{3}{m/s}, \num{1000}
\usepackage{relsize}         % \textlarger, \textsmaller for relative font sizing
\usepackage[version=4]{mhchem} % Chemical formulae: \ce{H2O}, \ce{CO2}, \ce{^1H-NMR}

% Floats
\usepackage{float}           % [H] = here definitely
\usepackage{stfloats}       % Improved placement for bottom floats
\usepackage{subcaption}     % Subfigures with (a), (b) and per-subfigure captions
\usepackage{rotating}       % \begin{sidewaystable}, \begin{sidewaysfigure} for wide content

% Algorithms
\usepackage{algorithm}       % Numbered algorithm floats and List of Algorithms
\usepackage[noend]{algpseudocode}

% Other useful packages
\usepackage{lipsum}          % Placeholder text: \lipsum[n], \lipsum[1-5], etc.
\usepackage{enumitem}       % Customise list labels, spacing
\usepackage{verbatim}       % \begin{verbatim} ... \end{verbatim} for code
\usepackage{afterpage}      % \afterpage{\clearpage} to clear after current page

% -----------------------------------------------------------------------------
% Dissertation metadata
% -----------------------------------------------------------------------------
\title{Template Dissertation Title: Title Case for Chapters}
\author{Your Full Name}
% Degree
\degree{BSc Computer Science}
% \degree{MSc Advanced Computer Science}
\university{Your University}
\date{February 2026}

% License and copyright following SPDX. (optional; comment out if not needed).
% Needs qdissertation-licenses.tex for license text in the License chapter.
\license{CC-BY-SA-4.0}
\copyrightholder{Your Full Name}
\licenseyear{2026}

\begin{document}

\frontmatter

% \maketitle{...} produces: title page, (copyright page if \license set), abstract,
% then frontlists: TOC, List of Figures, List of Tables, List of Algorithms,
% List of Acronyms (if acronyms.tex defines \newacronym).
% Add \dedication{...} and \makededication separately if desired.
\maketitle{%
  % Abstract: typically 200--300 words. State problem, methods, and main results.
  This is the abstract: a short, self-contained summary of the dissertation.
  It should state the problem, methods, and main results so that a reader
  can understand the scope without reading the full document. Replace with
  your own abstract.
}

% -----------------------------------------------------------------------------
% Optional front-matter chapters (unnumbered)
% -----------------------------------------------------------------------------

% Replace with your dedication (one or two lines); optional.
\dedication{%
  To my family.%
}

\chapter{Acknowledgements}
% Replace with your acknowledgements. Here you can also acknowledge any contributions from other
% parties to the work.
I thank my supervisors and colleagues for their support.

%\chapter{Declaration}
% Replace with your declaration (e.g. that the work is your own, and any
% required statements on collaboration, publications, or word count).
% Often this is submitted in separate forms or not needed, and does not
% have to be here.
%This dissertation is the result of my own work. Where reference is made to other work, this is indicated in the text.

%\chapter{Publications, data and code}
\chapter{Data and Code}
% List publications, if available or maybe planned, and state where data and code produced for this
% dissertation are available. Note, you usually must submit any sources and simlar you ahve
% written with your dissertation.
% Replace with your own entries; you can use \citep{key} for publications.

%\section{Publications}
% Papers (and optionally other outputs) that form part of or arise from this thesis.
% These can also be discussed in the introduction (what has been published or is planned).
%\begin{enumerate}
%  \item Author. Title. \textit{Journal}, volume, pages, year.~\citep{example-paper}
%\end{enumerate}

\section{Data}
% Datasets produced in this work and where they are available (repository, DOI, or ``on request'').
Data supporting this dissertation are available from [repository/DOI or ``available on request'']. Replace with your statement.

\section{Code}
% Software or code produced in this work and where it is available (repository, DOI, etc.).
Code produced for this dissertation is available from [repository/DOI or ``available on request'']. Replace with your statement.

\mainmatter

% -----------------------------------------------------------------------------
% Main chapters (suggested structure in each chapter file; adjust to your project)
% -----------------------------------------------------------------------------
\chapter{Introduction}\label{ch1}
% =============================================================================
% CHAPTER 1: Introduction
% =============================================================================
% SPDX-FileCopyrightText: 2026 Frank Langbein <frank@langbein.org>
% SPDX-License-Identifier: LPPL-1.3c
%
% Suggested sections: Motivation and context; Aims and objectives; Contributions;
% Dissertation structure (roadmap). Optional: Scope (boundaries of the work),
% Terminology or key terms. Start with a short chapter intro (no \section) before
% the first section. This file also demonstrates template features; see
% Section~\ref{sec:ch1-reference}. Remove that section for your own report.
% =============================================================================

% -----------------------------------------------------------------------------
% Chapter intro (no section): one short paragraph before the first \section
% -----------------------------------------------------------------------------
This chapter introduces the dissertation. It outlines the motivation and context for the work, the aims and objectives, the main contributions, and the structure of the remaining chapters. The following sections expand on each of these.

% -----------------------------------------------------------------------------
% Motivation and context (typical first section)
% -----------------------------------------------------------------------------
% Use \section, \subsection, \label{key}. Refer with \ref{key} or \autoref{key}.

\section{Motivation and context}\label{sec:motivation}

\newthought{This template illustrates} the structure and style of a dissertation using the Qyber dissertation class. As in a real introduction, you would state the problem domain, why it matters, and how your work fits. Acronyms use the long-short style: first use gives the full form (e.g.\@ \gls{uk}); later use \acrshort{api} or \acrlong{api} when you need only short or long.

As shown in previous work~\citep{example-paper}, we follow standard practice.
%F Clarify whether this applies to all cases.
\citet{example-book} provide a longer treatment; multiple citations: \citep{example-paper,example-conf}. You can link to other chapters (e.g.\@ Chapter~\ref{ch2}, Section~\ref{sec:objectives}) and to figures (Fig.~\ref{fig:demo}). The class provides \url{https://example.com} and \href{https://example.com}{link text} for URLs and hyperlinks.

% -----------------------------------------------------------------------------
% Aims and objectives (typical second section)
% -----------------------------------------------------------------------------
% Use \citep{} (parenthetical) or \citet{} (narrative), not \cite.

\section{Aims and objectives}\label{sec:objectives}

State the aim of your project and break it down into objectives. The following is a placeholder list (enumitem package: custom labels).

\begin{enumerate}[label=\textbf{O\arabic*:}, leftmargin=*]
  \item First objective: replace with your own.
  \item Second objective.
  \item Third objective.
\end{enumerate}

When you need mathematics, use amsmath and mathtools. Example:
\begin{equation}\label{eq:example}
  E = mc^2, \qquad \int_0^1 f(x)\,\mathrm{d}x = F(1) - F(0).
\end{equation}
Refer as Equation~\eqref{eq:example}; inline math: $\alpha + \beta = \gamma$. Paired delimiters (mathtools): \DeclarePairedDelimiter{\abs}{\lvert}{\rvert} then $\abs{x}$, $\abs*{\frac{1}{2}}$. Chemistry (mhchem): \ce{H2O}, \ce{CO2}, \ce{^1H-NMR}. Units (siunitx): \SI{3.0}{m/s}, \num{1234} samples. Symbols: \textdegree, \textcopyright; dingbats (pifont): \ding{51} \ding{43}. Relative size (relsize): \textlarger{larger} \textsmaller{smaller}.

% -----------------------------------------------------------------------------
% Contributions (typical third section)
% -----------------------------------------------------------------------------

\section{Contributions}\label{sec:contributions}

\epigraph{The best way to predict the future is to invent it.}{Alan Kay}

Summarise the main contributions of the dissertation. You can use a short list and refer to a figure or table. The class provides restyled \texttt{quote} and \texttt{quotation} environments (smaller type, indents):

\begin{quotation}
  A short quotation: replace with a real citation or excerpt. Use \verb|\begin{quotation}...\end{quotation}| for indented, smaller text. For a shorter block use \verb|\begin{quote}...\end{quote}|.
\end{quotation}

\begin{quote}
  A single-line quote: \verb|\begin{quote}...\end{quote}| gives a shorter block with smaller type (class restyle).
\end{quote}

The class also restyles \verb|\begin{verse}...\end{verse}| for short verse or poetry (centred lines, smaller type):
\begin{verse}
  Line one.\\
  Line two.\\
  Line three.
\end{verse}

\begin{itemize}
  \item First contribution; replace with your own.
  \item Second contribution.
\end{itemize}

Figure~\ref{fig:demo} shows how to include a figure (graphicx, caption; optional short caption for the List of Figures). Float placement: \texttt{[t]} = top of page (default here); \texttt{[htbp]} = here, top, bottom, float page; \texttt{[!t]} = relax limits and prefer top. See the ``Float placement'' comment block in \texttt{main.tex}. Use \texttt{[H]} (float package) only when you must force ``here'' (can cause uneven spacing). Subfigures (subcaption): see Figure~\ref{fig:sub}.

\begin{figure}[t]
  \centering
  \fbox{\rule{0pt}{4cm}\rule{4cm}{0pt}}
  \caption[Short for list of figures]{Full caption. Optional argument used in List of Figures.}\label{fig:demo}
\end{figure}

\begin{figure}[t]
  \centering
  \begin{subfigure}[t]{0.45\linewidth}
    \centering
    \fbox{\rule{0pt}{2cm}\rule{2cm}{0pt}}
    \subcaption{First subfigure.}
  \end{subfigure}
  \hfill
  \begin{subfigure}[t]{0.45\linewidth}
    \centering
    \fbox{\rule{0pt}{2cm}\rule{2cm}{0pt}}
    \subcaption{Second subfigure.}
  \end{subfigure}
  \caption{Example with subcaption (a) and (b).}
  \label{fig:sub}
\end{figure}

% Full-width figure: requires class option \texttt{fullwidth}. If you remove that option, comment out this block.
\begin{fullwidthfigure}[t]
  \fbox{\rule{0pt}{2.5cm}\rule{\linewidth}{0pt}}
  \caption[Full-width figure]{Full-width figure example. Use \texttt{\string
    \begin{fullwidthfigure}[t]}...\texttt{\string
  \end{fullwidthfigure}} when the class option \texttt{fullwidth} is set.}\label{fig:fullwidth}
\end{fullwidthfigure}

% Full-width block (not a float): extends into margin. Requires class option \texttt{fullwidth}; comment out if not set.
\begin{fullwidth}
  This paragraph is in \verb|\begin{fullwidth}...\end{fullwidth}|: block content that extends into the left and right margins (e.g.\@ for wide tables or key points).
\end{fullwidth}

Tables: use \texttt{booktabs} (\verb|\toprule|, \verb|\midrule|, \verb|\bottomrule|) (Table~\ref{tab:booktabs}), or \texttt{tabularray} with \texttt{tblr} / \texttt{longtable} for flexible or multi-page tables (Table~\ref{tab:demo}, Table~\ref{tab:long}). Same placement options as figures (\texttt{[t]}, \texttt{[!t]}, \texttt{[htbp]}, etc.).

\begin{table}[t]
  \centering
  \caption{Example table (booktabs, Tufte-style).}
  \label{tab:booktabs}
  \begin{tabular}{lcc}
    \toprule
    Item   & Value A & Value B \\
    \midrule
    First  & 1.0     & 2.0    \\
    Second & 3.0     & 4.0    \\
    \bottomrule
  \end{tabular}
\end{table}

\begin{table}[t]
  \centering
  \caption{Example table (tabularray).}
  \label{tab:demo}
  \begin{tblr}{colspec={lcc}, hlines}
    Item   & Value A & Value B \\
    First  & 1.0     & 2.0    \\
    Second & 3.0     & 4.0    \\
  \end{tblr}
\end{table}

\begin{center}
  \begin{longtable}{lcc}
    \caption{Minimal longtable demo (not in a float environment).} \label{tab:long} \\
    \toprule
    Item   & A & B \\ \midrule
    \endfirsthead
    \toprule
    Item   & A & B \\ \midrule
    \endhead
    Row 1  & 1 & 2 \\
    Row 2  & 3 & 4 \\ \bottomrule
  \end{longtable}
\end{center}

% -----------------------------------------------------------------------------
% Dissertation structure / roadmap (typical final section of introduction)
% -----------------------------------------------------------------------------

\section{Dissertation structure}\label{sec:roadmap}

Outline how the rest of the dissertation is organised so the reader knows what to expect. Reference chapters by label.

The rest of this dissertation is organised as follows.
\begin{itemize}
  \item \textbf{Chapter~\ref{ch2}} (Background): context, related work, problem statement.
  \item \textbf{Chapter~\ref{ch3}} (Methodology): research design, specification, AI/ML approach (if applicable), system design, implementation.
  \item \textbf{Chapter~\ref{ch4}} (Results and Evaluation): setup, results, evaluation, limitations.
  \item \textbf{Chapter~\ref{ch5}} (Conclusions and Future Work): summary and next steps.
  \item \textbf{Chapter~\ref{ch6}} (Reflection on Learning): transferable learning (omit if not required).
  \item \textbf{Appendix~\ref{app:supp}} and \textbf{Appendix~\ref{app:results}}: supplementary material and additional results.
\end{itemize}

Algorithms (algorithm, algpseudocode) appear in the List of Algorithms (Algorithm~\ref{alg:demo}). Use \texttt{[t]} or \texttt{[!t]} for top placement like figures/tables. For code snippets use verbatim; for wide tables or figures use \texttt{\string
\begin{sidewaystable}} or \texttt{\string
  \begin{sidewaysfigure}} (rotating).\footnote{Footnote example. Captions and footnotes are set at least 11\,pt by the class.} Placeholder prose in other chapters uses \texttt{\string\lipsum[n]}. Example: \lipsum[1]

    Use \texttt{\string\afterpage\{\string\clearpage\}} (afterpage) when you need to clear after the current page.

    \begin{algorithm}[t]
      \caption{Example algorithm (algpseudocode). Alternative: algorithmic for traditional layout.}
      \label{alg:demo}
      \begin{algorithmic}[1]
        \State \textbf{Input:} data $x$
        \State \textbf{Output:} result $y$
        \State $y \gets 0$
        \For{$i = 1$ \textbf{to} $n$}
        \State $y \gets y + x_i$
        \EndFor
        \State \Return $y$
      \end{algorithmic}
    \end{algorithm}

    Short code (verbatim):
\begin{verbatim}
  x = 1
  y = x + 1
\end{verbatim}

    % -----------------------------------------------------------------------------
    % Template reference (for authors: where to find package demos and conventions)
    % -----------------------------------------------------------------------------

    \section{Template reference}\label{sec:ch1-reference}

    This introduction doubles as a template: the sections above (Motivation, Aims and objectives, Contributions, Dissertation structure) are the usual structure for a dissertation Chapter~1. Replace their content with your own. The same file also demonstrates:

    \begin{itemize}
        \item \textbf{Citations and links:} \verb|\citep|/\verb|\citet| (natbib); link colour DarkBlue, cite/URL colour DarkGreen (class default).
      \item \textbf{\texttt{\string\newthought}:} Section-opening phrase in small caps with vertical space (e.g.\@ in Motivation). Letterspacing is applied by default.
        \item \textbf{Quote environments:} \verb|\begin{quotation}| \verb|...| \verb|\end{quotation}|, \verb|\begin{quote}| \verb|...| \verb|\end{quote}|, and \verb|\begin{verse}| \verb|...| \verb|\end{verse}| are restyled (smaller type, indents; verse centred) by the class.
        \item \textbf{Full-width content:} With class option \texttt{fullwidth}, use \verb|\begin{fullwidth}| \verb|...| \verb|\end{fullwidth}| for block content that extends into the margin, and \texttt{fullwidthfigure}/\texttt{fullwidthtable} for full-width floats (e.g.\@ Figure~\ref{fig:fullwidth}).
        \item \textbf{Epigraph:} \verb|\epigraph{quote}{source}| for chapter- or section-opening quotes (right-aligned, small type).
      \item \textbf{Other:} glossaries (long-short acronyms), amsmath/mathtools, textcomp, pifont, siunitx, relsize, mhchem, enumitem, graphicx/caption, float/stfloats, subcaption, tabularray/longtable, algorithm/algpseudocode, verbatim, afterpage, lipsum, rotating. Prefer parenthetical asides in the main text; use footnotes only sparingly (see \texttt{main.tex}).
    \end{itemize}

    Conventions for citations, spaces (\verb|~|, \verb|\@|), dashes, and reviewer comments (\verb|%ID |) are in the comment block at the top of \texttt{main.tex}. Class options (parskip/parindent, fullwidth, italic) are documented there too.


\chapter{Background}\label{ch2}
% =============================================================================
% CHAPTER 2: Background
% =============================================================================
% SPDX-FileCopyrightText: 2026 Frank Langbein <frank@langbein.org>
% SPDX-License-Identifier: LPPL-1.3c
%
% Suggested sections: Context (or Wider context / Domain and context); Related
% work (or Literature review / Prior work); Problem statement (and research
% question(s)). Optional: Theory and concepts; Stakeholders and constraints
% (as subsections of Context or separate). End with a clear problem statement.
% Assume a competent reader; avoid lengthy descriptions of standard tools.
% =============================================================================

This chapter provides the background needed to understand the project: the wider context, existing solutions or related work, and a clear problem statement (and research question(s) where applicable). Replace with your own review.

\lipsum[1-2]

\section{Context}\label{sec:context}
% Wider context, domain, motivation. Optional subsections: stakeholders,
% constraints, assumptions.

\lipsum[3-4]

\section{Related work}\label{sec:related}
% Prior solutions, methods, or systems; identify gaps your work addresses.
% Optional subsections: by theme, chronology, or type (e.g.\@ systems, methods, theory).

\lipsum[5-6]

\section{Problem statement}\label{sec:problem}
% Clear statement of the problem and, if applicable, research question(s).

\lipsum[7-8]


\chapter{Methodology}\label{ch3}
% =============================================================================
% CHAPTER 3: Methodology (or ``Specification, Design and Implementation'' / ``Approach'')
% =============================================================================
% SPDX-FileCopyrightText: 2026 Frank Langbein <frank@langbein.org>
% SPDX-License-Identifier: LPPL-1.3c
%
% This chapter describes how your solution works. Title it Methodology, Approach,
% or Specification, Design and Implementation; adapt or split as needed. Not all
% sections below apply to every project---pick, omit, or rename.
%
% Suggested sections by project type:
%   System design:  Specification, System design, Implementation. Optional: Datasets
%                   if data is central; Ethical considerations.
%   Research:       Research design, Literature and evidence (search strategy,
%                   inclusion criteria), Specification (= problem + requirements),
%                   Approach and solution, Implementation or experimental setup.
%   AI/ML:          Specification, Datasets, AI/ML models, System design (if a
%                   larger system), Implementation. Optional: Research design,
%                   Literature and evidence; Summary.
% =============================================================================

This chapter describes how the problem was addressed. Choose sections that fit your project type; omit or merge as needed. Replace with enough detail for reproduction.

\lipsum[9-10]

\section{Research design}\label{sec:research-design}
% Research: overall approach, methods (e.g.\@ quantitative, qualitative, design science), rationale. Omit for pure system-design projects.

\lipsum[11-12]

\section{Specification}\label{sec:spec}
% System design: requirements (what the system must do; optional subsections: functional, non-functional). Research: problem and what is required of a solution.

\lipsum[13-14]

\section{Datasets}\label{sec:datasets}
% AI/ML / empirical: data sources, preprocessing, train/validation/test split. Omit if not data-driven.

\lipsum[15-16]

\section{AI and machine learning models}\label{sec:aiml}
% AI/ML: model choice, training and validation, evaluation methodology. Omit otherwise.

\lipsum[17-18]

\section{System design}\label{sec:design}
% System design: architecture, data flow, behaviour (UML, ER, state charts). Use diagrams. Research: use ``Approach and solution'' for your proposed solution(s).

\lipsum[19-20]

\section{Implementation}\label{sec:impl}
% Realisation: key implementation choices and critical code; avoid describing every module. Optional: tools, platform, problems encountered.

\lipsum[21-22]

\section{Summary}\label{sec:summary}
% Optional: short summary of the methodology. Omit if the chapter is already concise.

\lipsum[23-24]

\chapter{Results and Evaluation}\label{ch4}
% =============================================================================
% CHAPTER 4: Results and Evaluation
% =============================================================================
% SPDX-FileCopyrightText: 2026 Frank Langbein <frank@langbein.org>
% SPDX-License-Identifier: LPPL-1.3c
%
% Suggested sections: Setup (or Evaluation setup / Experimental setup); Results;
% Discussion (or Evaluation and discussion); Limitations. Optional: Ablation
% studies (AI/ML); Comparison with baselines; Threats to validity. Describe how
% you demonstrated the system or hypothesis; include critical tests/experiments;
% critically appraise strengths and weaknesses.
% =============================================================================

This chapter presents the results and evaluates them. Describe how the solution was demonstrated, include test or experimental results, and critically appraise strengths and weaknesses.

\lipsum[25-26]

\section{Setup}\label{sec:setup}
% Evaluation setup: data, metrics, baselines, how experiments or tests were run.
% Optional subsections: environment, hyperparameters, evaluation protocol.

\lipsum[27-28]

\section{Results}\label{sec:results}
% Main findings: present tables and figures; summarise key outcomes. Optional
% subsections: by research question, by dataset, or by metric.

\lipsum[29-30]

\section{Discussion}\label{sec:discussion}
% Critical appraisal: strengths, weaknesses, methodology choices; comparison
% with prior work. Optional: ablation analysis, sensitivity analysis.

\lipsum[31-32]

\section{Limitations}\label{sec:limitations}
% Limitations, threats to validity, and suggested further tests if not all done.

\lipsum[33-34]


\chapter{Conclusions and Future Work}\label{ch5}
% =============================================================================
% CHAPTER 5: Conclusions and Future Work
% =============================================================================
% SPDX-FileCopyrightText: 2026 Frank Langbein <frank@langbein.org>
% SPDX-License-Identifier: LPPL-1.3c
%
% Suggested sections: Conclusions; Future work. Do not introduce new material in
% Conclusions---summarise aims and main results. Future work: next steps, open
% questions, what you would do with more time. Optional: short summary of
% contributions (as part of Conclusions or a bullet list).
% =============================================================================

This chapter concludes the dissertation. Summarise the aims, main results, and what was learned; then outline future work and improvements.

\lipsum[35-36]

\section{Conclusions}\label{sec:conclusions}
% Summary of aims, key results, and main takeaways. Optional: list of contributions.

\lipsum[37-38]

\section{Future work}\label{sec:future}
% Next steps, open questions, and what you would do with more time or resources.

\lipsum[39-40]


\chapter{Reflection on Learning}\label{ch6}
% =============================================================================
% CHAPTER 6: Reflection on Learning
% =============================================================================
% SPDX-FileCopyrightText: 2026 Frank Langbein <frank@langbein.org>
% SPDX-License-Identifier: LPPL-1.3c
%
% Suggested sections: Critical narrative of the experience; Analysis and application;
% Implications for future practice. Focus on transferable learning and (where
% applicable) ``double loop learning'': impact on assumptions and concepts, not
% only skills. Optional alternatives: ``What I would do differently''; ``Transferable
% skills and concepts''. Omit this chapter if not required by your institution;
% comment out the \chapter and \input in main.tex.
% =============================================================================

This chapter reflects on what was learned from the project: how assumptions and decisions were re-evaluated and what would be done differently in future.

\section{Critical narrative of the experience}\label{sec:challenges}
% Key challenges, turning points, and how they affected your understanding.

\lipsum[41-42]

\section{Analysis and application}\label{sec:reflection}
% How you analysed the experience and applied learning to decisions or concepts.

\lipsum[43-44]

\section{Implications for future practice}\label{sec:implications}
% What you will take forward: habits, assumptions, or approaches in future work.

\lipsum[45-46]


\backmatter

% -----------------------------------------------------------------------------
% Appendices
% -----------------------------------------------------------------------------
\appendix
\chapter{Supplementary Material}\label{app:supp}
% =============================================================================
% Appendix: Supplementary Material
% =============================================================================
% SPDX-FileCopyrightText: 2026 Frank Langbein <frank@langbein.org>
% SPDX-License-Identifier: LPPL-1.3c
%
% Appendices hold supporting material (detailed tables, derivations, survey text,
% user-study materials) so the main narrative stays focused. Complete source
% code is usually submitted separately, not in the report. Use \chapter and
% \section as in main chapters. Replace with your own supplementary material.
% =============================================================================

This appendix contains supplementary material: extra tables, derivations, or notation. Replace with your own content.

\section{Additional tables}\label{app:tables}

\lipsum[47-48]

\section{Further analysis}\label{app:analysis}

\lipsum[49-50]

\section{Derivation}\label{app:derivation}

\lipsum[51-52]


\chapter{Additional Results}\label{app:results}
% =============================================================================
% Appendix: Supplementary Material
% =============================================================================
% SPDX-FileCopyrightText: 2026 Frank Langbein <frank@langbein.org>
% SPDX-License-Identifier: LPPL-1.3c
%
% Appendices hold supporting material (detailed tables, derivations, survey text,
% user-study materials) so the main narrative stays focused. Complete source
% code is usually submitted separately, not in the report. Use \chapter and
% \section as in main chapters. Replace with your own supplementary material.
% =============================================================================

This appendix contains supplementary material: extra tables, derivations, or notation. Replace with your own content.

\section{Additional data}\label{app:data}

\lipsum[53-54]

\section{Further details}\label{app:details}

\lipsum[55-56]


% Bibliography (always last): \bib{bibfile} uses natbib (numeric or harvard per class option)
% and inserts the License chapter if \license was set.
\bib{bibliography}

\end{document}
